\chapter{System Implementation}
\label{ch:implementation}

This chapter describes how CampusLoop translates the architectural design of Chapter~\ref{ch:design} into concrete frontend, backend, persistence, worker, media, and deployment mechanisms. All module paths, function names, HTTP contracts, concurrency locks, and error codes reflect the verified codebase at \texttt{main@ecc22d2}.

\begin{clbox}{Implementation Pillars and Hardening Summary}
  \small
  \begin{itemize}[leftmargin=1.2em,itemsep=1pt,topsep=2pt]
    \item \textbf{End-to-end contract safety}: OpenAPI~3.1-generated TypeScript SDK paired with Pydantic~V2 \texttt{StrictInput}, uniform \texttt{BusinessError} envelopes, and signed 64-bit \texttt{BIGINT} bounds.
    \item \textbf{Deterministic concurrency}: five-tier row-lock ordering, \texttt{Idempotency-Key} payload hash verification, and transactional SQLAlchemy outbox hooks consumed via \texttt{SKIP LOCKED}.
    \item \textbf{Defense-in-depth media \& ops}: Pillow EXIF stripping and CRC validation into dual MinIO buckets, isolated audit databases, and public-only \texttt{HY-MT1.5-1.8B} catalog translation.
  \end{itemize}
\end{clbox}

\section{Frontend Architecture and Typed OpenAPI SDK Layer}
\label{sec:impl-frontend}

The web frontend (\texttt{frontend/src/}) is built with React and TypeScript \cite{reactDocs}, structured so that UI pages never assemble raw HTTP requests or hand-maintain response DTOs. Instead, the FastAPI backend exports its OpenAPI~3.1 specification \cite{openapi31}, from which \texttt{frontend/src/sdk/generated/sdk.gen.ts} and its accompanying TypeScript definitions are generated. A thin runtime wrapper in \texttt{frontend/src/sdk/index.ts} layers session management, automatic token refresh, idempotency headers, and error normalization over the generated client.

Key client-side reliability and security mechanisms implemented in \texttt{frontend/src/} include:
\begin{itemize}
  \item \textbf{In-memory identity restoration}: on initial page load, \texttt{restoreIdentity()} invokes \texttt{restoreBrowserSession} using the browser's \texttt{HttpOnly} refresh cookie together with the \texttt{X-CampusLoop-Browser} header, hydrating a short-lived access token strictly in JavaScript memory.
  \item \textbf{Scoped error and lock-out handling}: response interceptors in \texttt{sdk/index.ts} trigger local session cleanup on \texttt{401} or \texttt{423} only for authenticated or session-management requests (\texttt{sessionRequest || authenticationRequest}), preventing unauthenticated public fetches or third-party checks from spuriously clearing an active session.
  \item \textbf{Per-request cancellation and safe retry keys}: the \texttt{useRemote} hook attaches a fresh \texttt{AbortController} to every query so rapidly changing filters or route navigations cannot overwrite newer state with stale responses. Mutating forms use \texttt{useAttemptKey()} to bind a stable \texttt{Idempotency-Key} across retries of the same submission and explicitly call \texttt{reset()} once the operation succeeds so subsequent distinct edits obtain fresh keys.
  \item \textbf{Loop-free navigation}: \texttt{BackLink.tsx} sanitizes return paths by stripping URL fragments, rejecting external origins, preventing cycles, and capping navigation chain depth at eight entries.
\end{itemize}

\section{Backend API Layer, Validation, and Error Contracts}
\label{sec:impl-backend}

The backend (\texttt{backend/app/}) exposes modular FastAPI routers under \texttt{backend/app/\allowbreak api/\allowbreak routes/}: \texttt{auth.py}, \texttt{market.py}, \texttt{transactions.py}, \texttt{governance.py}, \texttt{moderation.py}, \texttt{admin\_\allowbreak console.py}, \texttt{intelligence.py}, and \texttt{realtime.py}. Every request body inherits from a shared Pydantic~V2 strict schema (\texttt{StrictInput}, configured with \texttt{extra="forbid"}) \cite{pydanticDocs} so unexpected client properties fail early with HTTP~\texttt{422} rather than being silently ignored.

All domain and validation errors funnel through \texttt{BusinessError} in \texttt{backend/app/\allowbreak core/\allowbreak errors.py}, returning a uniform JSON envelope containing a machine-readable error code (such as \texttt{AUTH\_\allowbreak ACCESS\_\allowbreak EXPIRED}, \texttt{MESSAGE\_\allowbreak ID\_\allowbreak CONFLICT}, \texttt{MEETUP\_\allowbreak VERSION\_\allowbreak CONFLICT}, \texttt{ORDER\_\allowbreak STATE\_\allowbreak CONFLICT}, or \texttt{INDEPENDENT\_\allowbreak REVIEW\_\allowbreak REQUIRED}) and human-readable message. Path-parameter lookups routed through \texttt{load()} enforce PostgreSQL signed 64-bit \texttt{BIGINT} range bounds ($1 \le \text{id} \le 2^{63}-1$) before querying the database, mapping out-of-range integers cleanly to \texttt{404 NOT\_FOUND} instead of raising unhandled database driver exceptions.

\section{End-to-End Source Chains}
\label{sec:impl-chains}

To provide direct code traceability across the full stack, Table~\ref{tab:source-chains} maps six core user and administrative workflows from React UI components through generated SDK operations and FastAPI route/service handlers to authoritative PostgreSQL tables and failure boundaries.

\begin{table}[htbp]
  \centering
  \footnotesize
  \setlength{\tabcolsep}{4pt}
  \renewcommand{\arraystretch}{1.12}
  \caption{Six end-to-end source chains verified at \texttt{main@ecc22d2}.}
  \label{tab:source-chains}
  \begin{tabularx}{\linewidth}{@{}>{\raggedright\arraybackslash}p{2.2cm}>{\raggedright\arraybackslash}p{4.7cm}>{\raggedright\arraybackslash}p{3.0cm}>{\raggedright\arraybackslash}X@{}}
    \toprule[1pt]
    \cltblhdr \textbf{\color{clNavy}Workflow} & \textbf{\color{clNavy}UI $\to$ SDK $\to$ Backend Handler} & \textbf{\color{clNavy}Persisted Tables} & \textbf{\color{clNavy}Primary Boundary / Guard} \\
    \midrule
    1. Browser Session & \texttt{AccountPages.tsx}, \texttt{sdk/index.ts} $\to$ \texttt{login}, \texttt{refreshSession}, \texttt{restoreBrowser\allowbreak Session} $\to$ \texttt{auth.py}, \texttt{services/auth.py} & \texttt{users}, \texttt{auth\_session\_\allowbreak families}, \texttt{refresh\_\allowbreak sessions}, \texttt{auth\_audits} & Rejects untrusted \texttt{Origin} or missing CSRF header; revokes full family on replay; expired token fails \texttt{401} before appeal check. \\
    2. Listing Publish & \texttt{MarketPages.tsx}, \texttt{ImageQueue.tsx} $\to$ \texttt{uploadImage}, \texttt{createProduct} $\to$ \texttt{market.py}, \texttt{services/storage.py} & \texttt{uploaded\_\allowbreak objects}, \texttt{products}, \texttt{product\_\allowbreak images}, \texttt{idempotency\_\allowbreak records} & Validates MIME/pixels/EXIF; enforces \texttt{owned\_keys} ownership \& purpose; conflicts on mutated idempotency replay. \\
    3. Chat \& Offer & \texttt{ChatPages.tsx} $\to$ \texttt{createChatSession}, \texttt{listMessages}, \texttt{sendMessage}, \texttt{createOffer} $\to$ \texttt{transactions.py} & \texttt{chat\_sessions}, \texttt{chat\_\allowbreak messages}, \texttt{offers}, \texttt{notifications} & Blocks self-contact; checks participant; enforces \texttt{MESSAGE\_ID\_\allowbreak CONFLICT}; dynamically marks past-expiry offers \texttt{EXPIRED}. \\
    4. Unique Reserve & \texttt{ChatPages.tsx} $\to$ \texttt{acceptOffer} $\to$ \texttt{transactions.py::\allowbreak accept\_offer} & \texttt{offers}, \texttt{products}, \texttt{orders}, \texttt{order\_events}, \texttt{idempotency\_\allowbreak records} & Blocks self-acceptance; locks \texttt{product -> offer}; Migration~\texttt{0003} unique \texttt{offer\_id} \& partial active \texttt{product\_id} index. \\
    5. Meetup \& Review & \texttt{TransactionPages.tsx} $\to$ \texttt{saveMeetup}, \texttt{confirmMeetup}, \texttt{confirmOrder\allowbreak Complete}, \texttt{createReview} $\to$ \texttt{transactions.py} & \texttt{meetups}, \texttt{orders}, \texttt{order\_events}, \texttt{products}, \texttt{reviews} & Rejects stale meetup version (\texttt{409}) and post-completion edits (\texttt{409}); dual flags for \texttt{COMPLETED}; allows review even if peer disabled. \\
    6. Governance \& Appeals & \texttt{AdminPages.tsx}, \texttt{AccountDecision\allowbreak Controls.tsx}, \texttt{AppealsPage.tsx}, \texttt{ModerationPages.tsx} $\to$ \texttt{resolveReport}, \texttt{overrideReport}, \texttt{adminCreate\allowbreak AccountDecision}, \texttt{submitAccount\allowbreak Appeal}, \texttt{adminReview\allowbreak AccountAppeal} $\to$ \texttt{governance.py}, \texttt{moderation.py} & \texttt{reports}, \texttt{account\_\allowbreak decisions}, \texttt{account\_\allowbreak appeals}, \texttt{auth\_audits}, \texttt{notifications} & Blocks self-reporting own message; separates triage from sanctions; revokes sessions \& protects orders on disable; enforces independent reviewer. \\
    \bottomrule
  \end{tabularx}
\end{table}

Several subtle operational refinements verified in \texttt{main@ecc22d2} close edge cases discovered during end-to-end auditing:
\begin{itemize}
  \item \textbf{Closing offers under state transitions}: \texttt{cancel\_offer} and \texttt{reject\_offer} pass \texttt{closing=True} when loading the target offer, enabling participants to cleanly close lingering \texttt{PENDING} offers even after the underlying product has been \texttt{RESERVED} by another trade or the counterparty's account has been disabled.
  \item \textbf{Dynamic expiry projection}: \texttt{offer\_fields()} dynamically reports any \texttt{PENDING} offer whose \texttt{expires\_at} timestamp has passed as \texttt{EXPIRED} during read serialization, preventing stale pending badges in chat even before a write touches the row.
  \item \textbf{Post-completion reviews with disabled peers}: \texttt{create\_review} invokes participant loading with \texttt{allow\_disabled\_peer=True} so an active student who already completed a legitimate transaction is not deprived of the ability to leave feedback if the counterparty is subsequently disabled.
  \item \textbf{Wanted-post expiry filtering}: \texttt{list\_wanted(status="EXPIRED")} includes both explicitly marked \texttt{EXPIRED} rows and \texttt{OPEN} rows whose \texttt{expires\_at} deadline has elapsed, ensuring filter views remain consistent with real clock time.
\end{itemize}

\section{Concurrency Control, Lock Ordering, and Idempotency}
\label{sec:impl-concurrency}

High-contention write paths in CampusLoop---such as two competing offers being accepted against the same listing, a buyer modifying a meetup while the seller confirms it, or an administrator disabling an account with in-flight orders---cannot rely on optimistic UI checks alone. The backend enforces three layered PostgreSQL concurrency guarantees \cite{postgresExplicitLocking}:
\begin{enumerate}
  \item \textbf{Deterministic five-tier row-lock hierarchy}: whenever a transaction locks multiple domain rows via \texttt{SELECT ... FOR UPDATE} (or \texttt{FOR NO KEY UPDATE} / \texttt{FOR SHARE NOWAIT} for read-integrity checks), it acquires locks strictly in the global order
  \[
    \texttt{auth / user} \;\longrightarrow\; \texttt{product} \;\longrightarrow\; \texttt{offer} \;\longrightarrow\; \texttt{order} \;\longrightarrow\; \texttt{meetup}.
  \]
  In particular, \texttt{services/order\_governance.py::protect\_disabled\_orders} sorts affected orders by \texttt{(product\_id, order\_id)} and locks \texttt{Product} before \texttt{Order} and \texttt{Meetup}, eliminating circular deadlock cycles with concurrent student trade mutations.
  \item \textbf{Database uniqueness backstops}: Migration~\texttt{0003} backs application locks with unique indexes on \texttt{orders.offer\_id}, active \texttt{orders.product\_id}, \texttt{meetups.order\_id}, \texttt{reviews(\allowbreak order\_id, reviewer\_id)}, and \texttt{chat\_\allowbreak messages(\allowbreak sender\_id, \allowbreak client\_\allowbreak msg\_id)}. Even if two requests execute across separate API workers, PostgreSQL enforces single-winner atomicity.
  \item \textbf{Request idempotency records}: critical mutating routes accept an \texttt{Idempotency-Key} header persisted in \texttt{idempotency\_records} alongside a canonical SHA-256 hash of the request payload and the serialized HTTP response. Replaying the exact request returns the cached result without duplicating rows; reusing the same key with a different payload fails with \texttt{409 IDEMPOTENCY\_CONFLICT}.
\end{enumerate}

\section{PostgreSQL Outbox Matching Worker and Observability}
\label{sec:impl-worker}

When a student publishes or edits a product, updates a wanted post, or changes account status, CampusLoop asynchronously recomputes matching recommendations for open wanted posts without slowing down the synchronous HTTP request. Rather than introducing a separate message broker that could drift out of sync with SQL transactions, \texttt{backend/app/services/matching\_jobs.py} implements a transactional PostgreSQL outbox pattern using SQLAlchemy session hooks \cite{sqlalchemy20}:
\begin{itemize}
  \item \textbf{Transactional event capture}: a \texttt{before\_flush} listener inspects dirty and new \texttt{Product}, \texttt{WantedPost}, and \texttt{User} instances across an explicit allowlist of business fields. In \texttt{after\_flush\_postexec}, the session increments the singleton \texttt{matching\_revision} counter, computes a deterministic SHA-256 \texttt{business\_event\_key} over normalized field values and timestamps, and executes an idempotent \texttt{INSERT ... ON CONFLICT DO NOTHING} into \texttt{matching\_jobs} for every active, non-expired \texttt{WantedPost}. If the enclosing transaction rolls back, \texttt{after\_rollback} clears staged changes so phantom jobs are never emitted.
  \item \textbf{Lock-held job claiming (\texttt{SKIP LOCKED})}: worker processes in \texttt{backend/app/\allowbreak worker/\allowbreak runner.py} claim pending rows via \texttt{SELECT ... FOR UPDATE SKIP LOCKED} \cite{postgresExplicitLocking} and hold the row lock until commit. A crashed worker automatically rolls back its transaction and releases the row lock without requiring lease-expiration sweeps. Failed jobs retry with bounded exponential backoff $\Delta t = \min(60, 2^{a})$ seconds up to \texttt{MAX\_ATTEMPTS = 5} before transitioning to \texttt{FAILED}.
  \item \textbf{Runtime observability}: the standalone worker attaches a \texttt{RunObserver} (\texttt{bp4-worker-\allowbreak runtime-v1}) backed by a bounded 2,000-event in-memory queue that records per-job execution telemetry in \texttt{ai\_task\_runs}, emits throttled heartbeats every 10 seconds to \texttt{worker\_heartbeats}, handles \texttt{SIGTERM}/\texttt{SIGINT} gracefully by draining in-flight commits, and serves live loop status on its internal \texttt{/healthz} HTTP endpoint.
\end{itemize}

\section{Media Security Pipeline and Dual-Bucket Isolation}
\label{sec:impl-media}

User-uploaded images are a common vector for polyglot payloads, decompression bombs, EXIF location leaks, and unauthorized object access. The upload pipeline in \texttt{backend/app/\allowbreak services/\allowbreak storage.py} enforces strict server-side sanitization and ownership isolation before persisting any binary object to MinIO \cite{minioDocs}:
\begin{enumerate}
  \item \textbf{Payload and pixel bounds}: raw uploads exceeding 5~MiB are rejected with \texttt{413 IMAGE\_TOO\_LARGE}. Pillow decodes the buffer in memory, allowing only \texttt{JPEG} and \texttt{PNG} formats whose total pixel count does not exceed $25\times 10^{6}$ pixels, and runs \texttt{image.verify()}.
  \item \textbf{Malicious/corrupted archive rejection}: malformed headers, truncated streams, decompression bombs, and corrupted PNG CRC chunks (\texttt{Unidentified\allowbreak ImageError}, \texttt{OSError}, \texttt{Decompression\allowbreak BombError}, \texttt{SyntaxError}, \texttt{ValueError}, \texttt{EOFError}, and \texttt{zlib.error}) are caught deterministically and mapped to \texttt{415 UNSUPPORTED\_IMAGE}.
  \item \textbf{EXIF stripping and JPEG re-encoding}: valid images are transposed according to EXIF orientation (\texttt{ImageOps.exif\_transpose}), converted to \texttt{RGB}, resized to $512\times 512$ via Lanczos resampling if \texttt{purpose == "avatar"}, and re-encoded into clean \texttt{JPEG} bytes (quality~90), stripping all original EXIF metadata and embedded chunks.
  \item \textbf{Dual-bucket separation and ownership verification}: product photos (\texttt{purpose == "product"}) are stored in \texttt{campusloop-public} under a UUID key; chat attachments, report evidence, decision evidence, and appeal evidence are stored in \texttt{campusloop-private} and assigned non-resolvable internal reference URIs (\texttt{https://private.campusloop.invalid/...}) so browsers can retrieve them only through authenticated FastAPI streaming routes. Every upload writes an \texttt{UploadedObject} row recording \texttt{(owner\_id, object\_key, purpose)}; attaching an image to a listing, message, report, decision, or appeal verifies ownership and purpose via \texttt{owned\_keys()}, rejecting cross-user or cross-purpose object substitution with \texttt{403 FORBIDDEN}. On database commit failure, the uploaded object is deleted in a compensating cleanup path. For appeal evidence specifically, \texttt{uploadAppealImage} verifies that the student currently holds an unrevoked, unappealed \texttt{AccountDecision}.
\end{enumerate}

\section{Runtime Profiles and Local Display Translation}
\label{sec:impl-ops}

To keep daily usage, automated testing, and acceptance audits strictly isolated, CampusLoop maintains separate database and Redis allocations. The local daily startup script (\texttt{scripts/start-local.ps1}) preserves persistent PostgreSQL and MinIO volumes across restarts, whereas write-heavy verification suites target dedicated isolated databases and Redis DB14/DB15 rather than the daily account store.

For bilingual classroom and cross-campus browsing, CampusLoop optionally pairs the web client with a standalone local translation microservice (\texttt{services/\allowbreak catalog\_translation\_\allowbreak server.py}) serving \texttt{HY-MT1.5-1.8B}. To preserve security and reliability boundaries, the translation server enforces four strict controls:
\begin{itemize}
  \item \textbf{CORS and MIME enforcement}: requests must originate from \texttt{ALLOWED\_ORIGINS} and carry \texttt{Content-Type: application/json} with a 10-second socket timeout.
  \item \textbf{Span preservation}: technical model codes, prices, and protected tokens are masked before generation and restored deterministically afterward so model names and currency figures are never hallucinated.
  \item \textbf{Per-item failure isolation}: batch catalog translations execute with \texttt{isolate\_\allowbreak failures=\allowbreak True} so a single problematic string falls back to its original text without aborting the entire listing page.
  \item \textbf{Strict public-text boundary}: only public product titles, descriptions, and search queries are routed to the translator; original database fields remain unchanged, and private chats, moderation facts, reports, and appeals are never transmitted to the translation service.
\end{itemize}
